\documentclass[10.5pt]{article}\usepackage[margin=.82in]{geometry}\usepackage[T1]{fontenc}\usepackage{lmodern,microtype,booktabs,longtable,amsmath,xcolor,hyperref,fancyhdr}
\definecolor{navy}{HTML}{173B57}\definecolor{pale}{HTML}{EAF3F6}\definecolor{good}{HTML}{E9F7ED}\hypersetup{colorlinks=true,linkcolor=navy}\pagestyle{fancy}\fancyhf{}\fancyhead[L]{\small Medical Assumption Memory Benchmark}\fancyhead[R]{\small v0.5.1}\fancyfoot[C]{\thepage}\setlength{\headheight}{14pt}
\title{\vspace{-1cm}\textcolor{navy}{\Huge Medical Assumption Memory Benchmark v0.5.1}\\\Large Matched-current-context dataset and verification protocol}\date{13 August 2026}\begin{document}\sloppy\emergencystretch=3em\maketitle
\noindent\colorbox{pale}{\parbox{.96\linewidth}{\textbf{Purpose.} Test whether an LLM can combine a current relational marker with raw history, an open-loop assumption, or verified cognitive memory to recover the oracle-defined support reaction. Correctness is computational, not medical.}}
\tableofcontents\newpage
\section{Research design}The three conditions are B0 raw longitudinal evidence, M0 open-loop prospective memory, and B1 verified relational cognitive memory. Every condition receives identical current observations, public reaction rules, response schema, and raw historical events. Only episode-specific memory organisation differs. Primary estimands are five-class macro-F1 differences B1-B0 and B1-M0. Supporting metrics include rule accuracy, grounded accuracy, rendering consistency, schema validity, citation validity, and latency.
\section{Dataset architecture}The dataset contains 40 reusable episode memories. Every memory is paired with all five reactions, giving 200 canonical contexts. Every canonical context has structured, chronological, and concise renderings, producing 600 contexts and 1,800 B0/M0/B1 calls. The split unit is memory: 8 development memories and 32 reserved memories. No memory or rendering crosses the split.
\begin{longtable}{p{.32\textwidth}p{.59\textwidth}}\toprule\textbf{Resource}&\textbf{Role}\\\midrule
\texttt{data\_v0\_5\_1/memories.csv}&40 prospective and verified relational memories.\\
\texttt{data\_v0\_5\_1/canonical\_contexts.csv}&200 contexts, 40 per reaction and 40 per template family.\\
\texttt{data\_v0\_5\_1/historical\_events.jsonl}&Source-level history and episode-specific comparison guide.\\
\texttt{data\_v0\_5\_1/current\_events.jsonl}&Balanced current observations and comparison marker.\\
\texttt{data\_v0\_5\_1/rendering\_map.csv}&Canonical/rendering/memory linkage.\\
\texttt{hidden/oracle\_references\_v0\_5\_1.jsonl}&Reaction, rule ID, hidden interpretation, evidence groups, hash.\\
\texttt{prompts/model\_calls\_v0\_5\_1.jsonl}&1,800 API-ready condition payloads.\\
\bottomrule\end{longtable}
\section{Matched current-context construction}Five current template families are used: recurrence, additional observation, new experience, source change, and temporal gap. A Latin-square allocation gives exactly eight examples of every reaction in every template family. Therefore template family is independent of reaction. Each memory also uses a random permutation mapping five neutral comparison markers to five relational meanings. The current text exposes one marker; current text alone cannot reveal its meaning. Raw history contains the episode-specific guide, and B1 preserves the same guide in compressed relational memory.

The five meanings are: comparable with no additional barrier; comparable with an active additional barrier; matched unresolved review boundary; different function; and missing comparison-critical field. These map to continuation, tailoring, escalation, withholding, and waiting under the disclosed oracle. Because the mapping is independently permuted for every memory, marker identity is not a global answer code.
\section{Cognitive-memory construction}The prospective record contains assumption, expected consequence, alternative, and formation evidence. Follow-up records describe implementation, observed target change, competing evidence, and timing limits. B1 preserves the resolved relation, provenance, scope, limitations, optional successor relation, and the episode-specific comparison guide. B1 does not contain the current marker selection or final reaction. Memory alone is therefore non-predictive; current evidence alone is also non-predictive. The answer requires their conjunction.
\section{Fairness and solvability}B0 may reconstruct the marker meaning from raw historical guide events. M0 receives the same history plus the prospective assumption. B1 receives the guide as structured memory. A deterministic B1 resolver must achieve 100\% oracle agreement by extracting the current marker, looking it up in the memory guide, and applying the public reaction rule. The oracle and renderer remain independent.
\section{Rendering design}All renderings preserve semantic event identity, dates, sources, comparison marker, evidence groups, oracle reference, and split. Rendering-oracle equivalence must be 100\%. Report macro-F1 by rendering family and all-render paraphrase consistency.
\section{External API contract}Every model response must be one JSON object containing reaction, rule ID, decisive event IDs, and a short reason. The local interpreter validates JSON and schema, reaction-rule consistency, visible citations, condition-aware historical and current evidence coverage, oracle correctness, and prohibited overreach. No second LLM or human adjudication is used.
\section{Challenge and integrity gates}Before external testing require: exact five-class balance; exact template-by-reaction balance; raw condition parity; correct B0/M0/B1 hierarchy; zero reaction-token leakage; deterministic B1 solvability; rendering equivalence; and grouped baseline thresholds. Current-only, assumption-only, and memory-only macro-F1 must be below 0.45; B0 bag-of-words below 0.55; identifier, structure, and event-count baselines near 0.20.
\section{Verification result}All structural and challenge checks passed. There are 40 memories, 200 canonical contexts, 600 renderings, and 1,800 calls. Every reaction and template has 40 canonical contexts; each template-reaction cell has eight. Development grouped baselines were: current-only 0.000 macro-F1, assumption-only 0.067, memory-only 0.067, and B0 bag-of-words 0.000. Identifier, structure, and event-count accuracy were each 0.20. The deterministic B1 resolver achieved 100\% oracle agreement. A complete 360-call development mock run achieved 100\% reaction, rule, grounding, and schema accuracy, confirming pipeline execution rather than model benefit.
\section{External execution sequence}\begin{enumerate}\item run structural validation and archive checksums;\item run retry, timeout, restart, and provider-structured-output QA;\item execute a 15-call development micro-smoke;\item execute the 360-call development set;\item freeze provider, model snapshot, prompts, parser, scorer, call order, retry policy, and seeds;\item execute 1,440 reserved calls;\item score deterministically and bootstrap by memory, retaining five contexts and all renderings per sampled memory.\end{enumerate}
\section{Interpretation boundary}Very low linear baselines show that shallow lexical models cannot learn the episode-specific marker permutation across held-out memories. This is intentional. The task remains deterministically solvable because the mapping is explicitly present in raw history and verified memory. External results test integration and retrieval of episode-specific relationships, not medical correctness.
\section{Go/no-go}\noindent\colorbox{good}{\parbox{.96\linewidth}{\textbf{GO for external API quality assurance and development testing.} Reserved execution remains conditional on production API retry, restart, and structured-output QA and a frozen release manifest.}}
\end{document}
